\documentclass[11pt,a4paper]{article}
\usepackage[T1]{fontenc}
\usepackage[utf8]{inputenc}
\usepackage{lmodern}
\usepackage[margin=25mm]{geometry}
\usepackage{amsmath,amssymb,amsthm}
\usepackage{microtype}
\usepackage{xcolor}
\usepackage{enumitem}
\usepackage{xurl}
\usepackage{bookmark}
\usepackage{hyperref}
\hypersetup{colorlinks=true,linkcolor=blue!45!black,urlcolor=blue!45!black, pdftitle={Bounded dependence and metric control: a dimension-one translation and a grid separation}, pdfauthor={AI-generated write-up}}
\setlength{\parindent}{0pt}
\setlength{\parskip}{0.45em}
\setlist{itemsep=0.25em,topsep=0.4em}
\setcounter{secnumdepth}{-1}
\setcounter{tocdepth}{1}
\providecommand{\tightlist}{\setlength{\itemsep}{0pt}\setlength{\parskip}{0pt}}
\emergencystretch=1.5em
\title{Bounded dependence and metric control\\[0.4em]\large A dimension-one translation and a grid separation}
\author{AI-generated write-up}
\date{12 October 2026}
\begin{document}
\maketitle
\begin{abstract}
A bounded-dependent outcome assigns a random output labeling to each input graph.
For locality $T$, the joint distribution on any vertex set depends only on its radius-$T$ input neighborhood, and output vectors on sets more than $2T$ apart are independent.
The outcome need not have a distributed sampling algorithm.

We show that, on bounded-degree graph classes of uniform Assouad--Nagata dimension at most one, such an outcome for a fixed locally checkable labeling problem gives a deterministic LOCAL algorithm in $O((T+\rho+1)\operatorname{polylog}n)$ rounds, provided its global success probability is at least $1-1/n$.
Here $\rho$ is the verification radius.
The algorithm constructs the required two-color cover and finite tables without querying the outcome's probabilities.

We also give an LCL on square grids without compass directions.
Its task is to assign edge bits whose XOR around every unit square is one.
It has a zero-error outcome with non-signaling locality at most three and dependence range two, but its deterministic LOCAL complexity is $\Theta(\sqrt n)$.
The non-signaling comparisons for this outcome are required only between grids.
Thus these restricted assumptions do not give a general polylogarithmic translation on grids.
The model sections state the size, identifier, and computability assumptions.
\end{abstract}
\clearpage
\tableofcontents
\clearpage
\part{Part I: Translation under one-dimensional metric control}
\section{1. Result and relation to prior work}\label{result-and-relation-to-prior-work}

We prove the following translation.
On a bounded-degree graph class of uniform Assouad--Nagata dimension at most one, a bounded-dependent outcome of locality \(T\) solving an LCL with global probability at least \(1-1/n\) yields a deterministic LOCAL algorithm with round complexity

\[ O\bigl((T+\rho+1)\operatorname{polylog} n\bigr), \]

where \(\rho\) is the verification radius.
For a fixed LCL and \(T\ge 1\), this is \(O(T\operatorname{polylog} n)\).
The algorithm need not sample the outcome or query its probabilities.
The metric coloring is assumed to exist and is constructed by the algorithm.

More generally, suppose every graph admits, at every integer scale \(r\ge 1\), a two-coloring whose monochromatic \(r\)-components have diameter at most \(Cr+A\), for common known computable integers \(C\ge 1\) and \(A\ge 0\).
Our bound is

\[ (T+A+\rho+1)(2+\log_2 N)^{O(\log(C+2))}. \]

Here \(N\) is the common size information: either the exact size \(n\), or a polynomial upper bound \(n\le N\le n^c\) for fixed \(c\).
The exponent has an absolute implicit constant.
Fixed \(C,A\) give the first statement.

\href{https://arxiv.org/html/2504.05191v2\#S4}{Balliu et al., Definitions 4.2--4.4} define the bounded-dependent outcome model used here.
Their \href{https://arxiv.org/html/2504.05191v2\#S4.SS1}{Corollary 4.13} gives a general deterministic translation with complexity \(\widetilde O(\sqrt{nT})\) and success strictly greater than \(1-1/n\).
We use their partial-sampling principle, but obtain the improved class-dependent bound from a two-color metric decomposition.
Counting failures by component also permits success exactly \(1-1/n\).

The standard coloring formulation of Assouad--Nagata control is given in \href{https://arxiv.org/html/2308.10377v3\#S1.SS4}{Distel, Section 1.4 and Proposition 5}.
His \href{https://arxiv.org/html/2308.10377v3\#S3}{Theorem 13} supplies linear two-color control for each fixed treewidth bound.
\href{https://arxiv.org/html/2308.12273v5\#S1.SS1}{Liu independently proves the minor-closed dimension bound}: proper minor-closed weighted graph classes have Assouad--Nagata dimension at most two, and his \href{https://arxiv.org/html/2308.12273v5\#S1.SS2.SSS4}{Theorem 1.7} gives dimension at most one for fixed treewidth.
Our translation uses dimension-one control; the general dimension-two theorem does not by itself provide the two colors needed here.
Our construction applies to any class with the stated two-color control.
\href{https://arxiv.org/abs/2608.16854v1}{d'Amore and Lievonen} obtain a stronger, tree-specific conclusion: subpolynomial bounded dependence gives an \(O(\log n)\) deterministic algorithm on trees.

The algorithmic input from prior work is deterministic network decomposition.
\href{https://arxiv.org/abs/1907.10937}{Rozhoň and Ghaffari} established polylogarithmic deterministic network decomposition; we use the parameters of \href{https://arxiv.org/html/2410.19516v1\#S1}{Ghaffari and Grunau, Theorem 1.1}: \(O(\log n)\) colors, \(O(\log n)\) diameter, and \(\widetilde O(\log^2 n)\) construction time.
All other ingredients needed for the translation are proved below.

\section{2. Model and hypotheses}\label{model-and-hypotheses}

\subsection{2.1. Inputs and finite computation}\label{inputs-and-finite-computation}

Graphs are finite, nonempty, simple, undirected, and have maximum degree at most a fixed constant \(\Delta\).
Vertices have unique IDs from a specified finite range of size polynomial in \(N\), local port numbers, and input labels from a finite alphabet.
Outputs also belong to a finite alphabet \(\Sigma\).
Edge and half-edge outputs can be encoded in vertex-output tuples ordered by ports; bounded degree keeps this alphabet finite.

The LOCAL model has synchronous rounds, unbounded message length, and unbounded local memory and computation.
Every local computation must nevertheless terminate.
The theorem bounds communication rounds, not sequential computation time.

Let \(\mathcal I_N\) be the allowed input instances with common size information \(N\).
We use either of these conventions:

\begin{enumerate}
\def\labelenumi{\arabic{enumi}.}
\tightlist
\item \textbf{Exact size:} \(N=n\), and \(\mathcal I_N\) consists of allowed \(N\)-vertex instances.
\item \textbf{Supplied upper bound:} \(n\le N\le n^c\), where \(c\) is fixed.
The exact \(n\) is not separately supplied, and \(\mathcal I_N\) contains all allowed instances compatible with that same \(N\).
\end{enumerate}

A terminating program must produce the complete finite list \(\mathcal I_N\) from \(N\) and the fixed problem description.
Decidable membership in the promise family suffices: enumerate all candidate graphs of at most \(N\) vertices, with the prescribed finite IDs, ports, and inputs, and test them.
A merely enumerable promise is insufficient because subsequent searches must know when they have tested every instance.
All algorithmic bounds, including \(T,k,C,A\), are known computable integers or functions of \(N\).

An LCL \(\Pi\) has verification radius \(\rho\) and finite input and output alphabets.
Write \(\Pi(G)\) for its valid global outputs, with the input data implicit.
A verifier's scope is its whole radius-\(\rho\) vertex ball, including coordinates it may ignore.
A scope is nonempty and has diameter at most \(2\rho\).
All admissible one-vertex instances are assumed solvable: the success threshold \(1-1/n\) alone says nothing when \(n=1\).

The outcome, the actual \(n\), and the supplied \(N\) always refer to the original whole instance.
Solving a connected component does not create a new outcome instance or change the size parameter.
Distances between different connected components are infinite.

\subsection{2.2. Whole-set locality and independence}\label{whole-set-locality-and-independence}

For \(U\subseteq V(G)\), its \textbf{marked radius-\(T\) input view} is the union of the information visible within \(T\) communication rounds of vertices of \(U\), with marks identifying the output coordinates in \(U\).
This view includes IDs, ports, degrees, input labels, and supplied size information, with the usual boundary-edge convention.
An extra gathering round can supply boundary information if needed.
Isomorphisms preserve all these data and identify marked output coordinates.

For every \(G\in\mathcal I_N\), suppose there is a distribution \(\mu_G\) on \(\Sigma^{V(G)}\).
Invalid outputs are allowed.
Let \(\mu_{G,U}\) denote its joint marginal on \(U\).
We require:

\begin{enumerate}
\def\labelenumi{\arabic{enumi}.}
\item \textbf{Whole-set non-signaling locality:} \(\mu_{G,U}\), including all its probabilities, is determined by the marked radius-\(T\) input view of \(U\).
Isomorphic marked views in \(\mathcal I_N\) have matching joint marginals.
\item \textbf{Whole-set bounded dependence:} for arbitrary vertex sets \(U,V\),

\[ d_G(U,V)>k\quad\Longrightarrow\quad \mu_{G,U\cup V}=\mu_{G,U}\otimes\mu_{G,V}. \]
\item \textbf{Global success:} for \(n=|V(G)|\),

\[ \mu_G(\Pi(G))\ge 1-\frac1n. \]
\end{enumerate}

The main specialization is \(k=2T\), which is a bounded-dependent outcome of locality \(T\) in Balliu et al.'s terminology.
If ``\(T\)-dependent'' refers instead to independence beyond distance \(T\), the same theorem applies with \(k=T\) and locality radius at most \(T\).

The sets in both locality conditions may be disconnected.
In particular, independence is required between a component and a union of other components.
Pairwise independence of individual vertex outputs, or even pairwise independence of component vectors, does not by itself imply the mutual independence used below.
A per-verifier success guarantee does not replace global success.
The probabilities may be irrational and noncomputable; no probability oracle is assumed.

A \(T\)-round randomized LOCAL algorithm with independent private random tapes satisfies the two locality conditions with \(k=2T\): outputs on \(U\) depend only on its radius-\(T\) input view and the tapes there, and radius-\(T\) neighborhoods of sets at distance greater than \(2T\) are disjoint.
This explains one source of outcomes satisfying the theorem, but the theorem does not assume the outcome can be sampled by a LOCAL algorithm.

\subsection{2.3. Uniform metric control}\label{uniform-metric-control}

For a set \(X\), its weak diameter is

\[ \operatorname{wdiam}_G(X)=\max_{x,y\in X}d_G(x,y). \]

All distances are in the original graph \(G\); shortest paths may leave \(X\).
The \textbf{\(r\)-components} of \(X\) are the connected components of the graph on \(X\) joining vertices at ambient distance at most \(r\).
Thus an \(r\)-chain may use ambient paths through vertices outside \(X\).

Assume that for each integer \(r\ge 1\), every allowed graph has a red--blue coloring whose monochromatic \(r\)-components have weak diameter at most

\[ f(r)\le Cr+A. \]

The constants are uniform across the class.
Uniform Assouad--Nagata dimension at most one supplies this property with \(A=0\) and fixed \(C\), after enlarging \(C\) to an integer.
The affine version also permits fixed additive control.
Requiring a graph-dependent control constant separately for each finite graph would not give this theorem.

Restricting a witnessing coloring to a vertex subset preserves its control in the restricted ambient metric: every monochromatic chain in the subset is a chain in the original graph.
We will use this fact on subsets which need not themselves be admissible LCL instances.

\section{3. Constructing the structural two-coloring}\label{constructing-the-structural-two-coloring}

This section uses only metric control and network decomposition.
The red and blue colors organize the algorithm and are unrelated to the LCL output alphabet.

\subsection{3.1. Discovering weak-diameter components}\label{discovering-weak-diameter-components}

\textbf{Lemma 1.} Suppose membership in a marked set is locally known, and its \(a\)-components have known weak-diameter bound \(D\).
In \(O(D+a+1)\) rounds, every marked vertex can find its whole \(a\)-component and its minimum-ID leader.

\textbf{Proof.} The entire component of a vertex \(v\) lies in \(B_G(v,D)\).
Gather the marked radius-\(3(D+a)+1\) neighborhood.
This determines exact ambient distances among vertices of \(B_G(v,D)\): two such vertices have a path through \(v\) of length at most \(2D\), so a shortest path lies within radius \(3D\) of \(v\).
Construct the \(a\)-graph on marked vertices in \(B_G(v,D)\) and find the component of \(v\).
Every true component vertex and every true component step is present.
Conversely, each computed step is a true ambient step.
The computed component is exact.
Choose its minimum-ID vertex.
All gathers can run concurrently with unbounded messages. \(\square\)

The same observation lets a leader gather a region of weak diameter \(D\), its fixed-radius neighborhood, and the ambient distance matrix required for a finite search, in time proportional to these radii.
Induced connectivity is unnecessary.

\subsection{3.2. Extending over separated regions}\label{extending-over-separated-regions}

\textbf{Lemma 2.} Let \(Z\subseteq V(G)\) have a partial red--blue coloring whose monochromatic \(t\)-components have weak diameter at most \(R\), with integer \(t\ge1\) and \(R\ge0\).
Let \(K_i\subseteq Z\) be disjoint uncolored regions, with locally known region IDs and designated leaders, weak diameter at most \(D\), and

\[ d_G(K_i,K_j)>R+2t\qquad(i\ne j). \]

All regions can be colored concurrently, preserving prescribed colors in \(Z\), in \(O(D+R+t+1)\) rounds.
The resulting monochromatic \(t\)-components have diameter at most

\[ F_t(R)=(C+2)R+(2C+2)t+C+A. \]

\textbf{Proof.} First consider a finite metric subspace \(X\) with prescribed colors on \(P\subseteq X\), whose monochromatic \(t\)-components have diameter at most \(R\).
Take an existing two-coloring of \(X\) at scale \(R+2t+1\) and overwrite its colors on \(P\) by their prescribed values.
Call the unchanged vertices in \(X\setminus P\) anchors.

For a final red \(t\)-chain, consecutive anchors are at distance at most \(R+2t\): the prescribed red vertices between them belong to one old red \(t\)-component.
If no prescribed vertex is between them, their distance is at most \(t\).
Hence all anchors of a final red component lie in one base red \((R+2t+1)\)-component, whose diameter is at most \(f(R+2t+1)\).
Every non-anchor in a component with an anchor is at distance at most \(R+t\) from an anchor, by following the chain to its first anchor.
The component diameter is therefore at most

\[ f(R+2t+1)+2(R+t)\le F_t(R). \]

A component with no anchor lies in an old component and has diameter at most \(R\le F_t(R)\).
The argument for blue is identical.
This proves existence of an extension on \(X\).

For a current region \(K\), let \(P_K\) be the union of all complete old monochromatic \(t\)-components within distance \(t\) of \(K\), of either color.
If \(v\) leads \(K\), then

\[ K\cup P_K\subseteq B_G(v,D+R+t). \]

The leader gathers radius \(3(D+R+t)+1\).
Region marks identify \(K\); working colors and the old diameter bound identify the complete components in \(P_K\); the gathered graph determines ambient distances in \(K\cup P_K\).
The restriction of metric control to this subspace proves that some coloring of \(K\) extends \(P_K\) with bound \(F_t(R)\).
The leader enumerates all two-colorings of \(K\) and selects the first passing the component-diameter test.
The search is finite and has a solution.

An old component cannot touch two current regions within distance \(t\): its diameter would imply their distance is at most \(t+R+t\).
There is also no direct \(t\)-step between regions.
A new monochromatic chain cannot consequently connect two current regions through old vertices.
Every changed component lies in some \(K\cup P_K\), where its diameter was tested, and untouched components retain their old bound.
The local choices are compatible globally.
Gathering and distributing costs \(O(D+R+t+1)\) rounds.
Colors outside \(Z\) are not prescribed coordinates of this operation; all communication still uses \(G\). \(\square\)

\subsection{3.3. Halving a palette}\label{halving-a-palette}

Set

\[ B=4C+8,\qquad M=B^4. \]

If \(R\ge t\), \(R\ge A\), and \(R\ge1\), then

\[ F_t(R)\le(4C+5)R\le BR. \]

\textbf{Lemma 3.} Suppose \(q\ge4\) is divisible by four, and a \(q\)-coloring has monochromatic \(a\)-components of weak diameter at most \(D\), where \(a=Mt\) and \(t\ge A+1\).
In \(O(D+a)\) rounds it can be replaced by a \(q/2\)-coloring whose monochromatic \(t\)-components have diameter at most \(a\).

\textbf{Proof.} Partition the old color names into groups of four.
For one group, let \(Z\) be its vertices.
Each vertex saves its old color and old monochromatic \(a\)-component ID, using Lemma 1.
Erase the working colors on \(Z\).
Process its four old color classes in order, treating their saved components as the current regions in Lemma 2.

Use \(R_i=B^it\) as the claimed component-diameter bound after \(i\) phases, with \(R_0=t\) for the empty coloring.
The preceding inequality gives \(R_{i+1}\) after the next phase.
Before each phase \(R_i\le B^3t\).
Distinct saved components of the current old color are at distance greater than \(a\), and

\[ a=B^4t>B^3t+2t, \]

since \(B\ge12\).
Lemma 2 therefore applies.
After four phases the bound is \(B^4t=a\).

Run all groups concurrently, giving each group its own pair of new color names.
Their vertex sets are disjoint, messages can be bundled, and disjoint new palettes prevent a monochromatic component from crossing groups.
Component discovery and each of four phases cost \(O(D+a)\), with no factor for the number of groups. \(\square\)

\subsection{3.4. From network decomposition to two colors}\label{from-network-decomposition-to-two-colors}

Fix computable network-decomposition guarantees \(Q=O(2+\log N)\) colors, \(d=O(2+\log N)\) weak diameter, and \(P=(2+\log N)^{O(1)}\) rounds, available on arbitrary graphs.
The standard construction can use the supplied polynomial size bound and ID range.
Pad \(Q\) to a power of two with \(Q\ge4\).
The guarantees do not depend on the maximum degree of the graph on which decomposition is run.

For a positive integer \(r\), one communication round on the power graph \(G^r\) can be simulated in \(O(r)\) rounds on \(G\) by relaying messages, including the initial discovery of power-graph neighbors and ports.
A decomposition cluster of weak diameter \(d\) in \(G^r\) has weak diameter at most \(rd\) in \(G\).
Distinct clusters of the same decomposition color are at distance greater than \(r\) in \(G\).

\textbf{Lemma 4.} For any integer \(s\ge A+1\), a deterministic LOCAL algorithm constructs a two-coloring whose monochromatic \(s\)-components have diameter at most \(D=Ms\), in

\[ s(2+\log_2N)^{O(\log(C+2))} \]

rounds, without a supplied witnessing coloring.

\textbf{Proof.} Put

\[ h=\log_2(Q/2),\qquad r=M^hs. \]

Compute a network decomposition of \(G^r\) and color vertices by cluster color.
Every monochromatic \(r\)-component is contained in one cluster, so its diameter is at most \(D_0=rd\).

Apply Lemma 3 for \(h\) passes.
After pass \(j\), there are \(Q/2^j\) color names, the component scale is \(r/M^j\), and, for \(j\ge1\), the diameter bound is \(r/M^{j-1}\).
Each target scale is at least \(s\ge A+1\).
The final palette has two colors, scale \(s\), and diameter \(Ms\).

The first pass costs \(O(r(d+1))\).
Subsequent costs form a geometric sum \(O(r)\), because both input diameter and scale decrease geometrically.
Including the initial decomposition, the total is

\[ O\bigl(r(P+d+1)\bigr). \]

Finally,

\[ r=s(Q/2)^{\log_2M},\qquad \log_2M=4\log_2(4C+8)=O(\log(C+2)). \]

Substitution gives the bound.
Every stage has a common schedule determined by known bounds; empty colors simply use their allotted time.
No global termination detection is needed.
Metric control is applied only in the original ambient metric, never to \(G^r\). \(\square\)

\section{4. Independent red choices and component obstructions}\label{independent-red-choices-and-component-obstructions}

Choose

\[ s=1+\max\{k,2\rho,A,1\},\qquad D=Ms, \]

and construct the coloring from Lemma 4.
Thus \(s\ge2\) and \(s>k,2\rho\).

Solve each ordinary connected component with at most two vertices directly.
It can gather its entire input and search for a valid labeling in \(O(\rho+1)\) rounds.
For \(n\ge2\), global success is positive and hence such a labeling exists; for \(n=1\), use the explicit solvability assumption.
Choose the first solution in ID and alphabet order.
No verifier crosses between ordinary connected components.
Keep the original \(G,n,N\) and omit these solved components from the following objective.

Let \(U_1,\ldots,U_m\) be the remaining red \(s\)-components.
Introduce independent random variables \(Z_i\) taking output patterns on \(U_i\), with law \(\mu_{G,U_i}\).
These variables are whole-component patterns, not independent choices for individual vertices.

Their product law agrees with the original red marginal:

\[ \mu_{G,\bigcup_iU_i}=\bigotimes_i\mu_{G,U_i}. \]

Indeed, \(U_j\) is at distance greater than \(s>k\) from the union of all earlier components.
Apply whole-set independence to those two sets and induct on \(j\).
This proves mutual independence, rather than assuming it from pairwise component independence.

Assign every remaining verifier whose scope contains blue vertices to the unique blue \(s\)-component containing those vertices.
Assign every red-only verifier to the unique red \(s\)-component containing its scope.
These assignments are well defined because scope diameter is at most \(2\rho<s\).

For each remaining structural component \(K\), define a zero--one obstruction \(E_K\):

\begin{itemize}
\tightlist
\item For red \(K\), it is one if some assigned red-only verifier rejects its chosen pattern.
\item For blue \(K\), it is one if no labeling of \(K\) satisfies all assigned verifiers with the red labels fixed.
\end{itemize}

If no verifier is assigned to \(K\), set \(E_K=0\).
Blue feasibility is a finite search over \(\Sigma^K\).
Let

\[ F=\sum_KE_K,\qquad b=\#\{\text{remaining structural components}\}. \]

Then \(F=0\) precisely when all chosen red labels admit a valid completion.
If every blue search succeeds, combine the passing patterns: no verifier reads blue variables from two components, and red-only checks pass.
Conversely, any valid completion supplies passing patterns for all blue searches and satisfies all red-only checks.

\textbf{Lemma 5.} The true product law of the red variables satisfies

\[ \mathbb E[F]\le1-\frac1N<1 \]

when \(N\ge2\).

\textbf{Proof.} For each \(K\), a valid original sample makes \(E_K=0\) on its red restriction: its blue labels witness feasibility if \(K\) is blue.
The law of that red restriction is our product law.
Therefore

\[ \Pr(E_K=1)\le\frac1n, \qquad \mathbb E[F]\le\frac bn. \]

No independence between the indicators is used.

If remaining vertices exist, at least one remaining ordinary connected component has at least three vertices, hence contains a path of three distinct vertices.
Its three vertices have pairwise distances at most two.
Two have the same structural color; since \(s\ge2\), those two lie in one monochromatic \(s\)-component.
Thus at least one part of the structural partition contains two vertices.
If \(n'\) vertices remain, \(b\le n'-1\le n-1\).
Consequently

\[ \mathbb E[F]\le\frac{n-1}{n}=1-\frac1n\le1-\frac1N. \]

If no vertices remain, \(F=0\) and the bound also holds.
The exact-size case \(N=1\) was already solved. \(\square\)

Grouping by component is essential to this endpoint argument.
Counting one obstruction per verifier could give \(n\) terms, leaving only \(\mathbb E[F]\le1\) at success exactly \(1-1/n\).

\section{5. Computing usable probability tables without an oracle}\label{computing-usable-probability-tables-without-an-oracle}

The true red marginals prove that suitable probability tables exist; they need not be computable.
We now synthesize different rational tables, using only a finite search.

\textbf{Lemma 6.} For \(N\ge2\), every vertex can compute the same finite table assigning a dyadic probability vector to each marked radius-\(T\) input type of a remaining red component, such that the resulting independent red choices satisfy

\[ \mathbb E[F]\le1-\frac1{2N} \]

on every instance in \(\mathcal I_N\).

\textbf{Proof.} Generate the complete instance list and simulate the deterministic structural construction and preprocessing on each instance.
List all remaining red components and their marked radius-\(T\) input types.
Structural colors select which sets occur; they are not extra marks in the input type used by non-signaling locality.
Canonicalize each finite type, with its marked vertices ordered by ID, and list all patterns on those vertices.
There are finitely many types and patterns.

For each integer \(a\ge1\), enumerate all tables whose entries are nonnegative multiples of \(2^{-a}\) and sum to one for each type.
A common denominator suffices for every type.
Test each table on every listed input.
For an input, enumerate all tuples of component patterns, compute the integer \(F\) by the local-verifier and blue-feasibility tests, and calculate its expectation using the candidate product probabilities.
This is a finite rational computation.
Accept the first table satisfying the displayed inequality for every instance.
Inputs already solved in preprocessing have \(F=0\).

To prove termination, whole-set locality supplies a well-defined true probability vector for each type.
Lemma 5 gives a margin of at least \(1/N\) below one on every instance.
Each tested expectation has the form

\[ \sum_{z_1,\ldots,z_m} F(z_1,\ldots,z_m)\prod_{i=1}^mp_{\tau_i}(z_i), \]

where \(\tau_i\) is the component's type.
This is a continuous polynomial in finitely many table entries, including when types repeat.
Since there are finitely many instances, sufficiently accurate simultaneous approximations of all true vectors change every expectation by less than \(1/(2N)\).

Dyadic probability vectors can approximate each finite vector arbitrarily closely: round down all but one coordinate at a sufficiently large denominator and put the remaining mass in the last coordinate.
The last coordinate is nonnegative and all coordinate errors tend to zero.
Thus some enumerated table passes.
Every candidate test terminates, and a passing candidate exists, so the search terminates.
It uses no true probability values, no unknown actual graph, and no communication.

All vertices use the same parameters, finite instance list, and enumeration order.
A red leader gathers its marked radius-\(T\) input view and looks up its table entry.
Matching occurrences have the same coordinates under view isomorphisms. \(\square\)

The synthesized tables need not reproduce the original marginals or extend to a global non-signaling law.
Only their product law and tested expectation bound are used.
Repeated occurrences of one type use the same probability vector but are independent random variables.
This is a single uniform algorithm across sizes, not size-dependent uncomputable advice.

\section{6. Deterministically fixing the red variables}\label{deterministically-fixing-the-red-variables}

\textbf{Lemma 7 (local conditional expectation).} Suppose finite independent variables with rational probability tables are located at vertices of an \(n\)-vertex graph.
Zero--one functions \(E_v\) are located at designated vertices.
For each \(E_v\), all its variable locations, function description, and probability tables can be learned within radius \(L\ge1\).
If \(\mathbb E[\sum_vE_v]<1\), then every variable can be fixed deterministically so that all \(E_v=0\) in

\[ O\bigl(L[P+Q(d+1)]\bigr)=L\operatorname{polylog}N \]

rounds.

\textbf{Proof.} Compute a network decomposition of \(G^{2L}\) in \(O(LP)\) physical rounds.
Clusters have weak diameter at most \(2Ld\) in \(G\), and distinct clusters with one color are at distance greater than \(2L\).

Process decomposition colors sequentially.
In a current cluster \(W\), the minimum-ID leader gathers its variables and every indicator depending on one of them, with that indicator's function description, all its needed probability tables, and already fixed values.
An affected indicator is within \(L\) of a variable in \(W\), and all its variable locations are within a further \(L\).
This data can be gathered and results distributed in \(O(L(d+1))\) rounds, using ambient paths even outside \(W\).

The leader enumerates positive-probability assignments to the unfixed variables in \(W\).
For each candidate, compute exactly the conditional expectation of the sum of affected indicators, holding prior choices fixed.
The weighted average over candidates equals its value before fixing \(W\).
Some candidate does not increase this expectation; choose the first such candidate and distribute its values.
All comparisons are finite sums of rational numbers.

Two current clusters cannot affect the same indicator: such an indicator would have variable locations in both clusters within distance \(L\) of its center, contradicting the clusters' separation greater than \(2L\).
Furthermore, fixing a cluster does not alter expectations of unaffected indicators, since unassigned coordinates retain their product law.
Adding the inequalities for concurrently acting clusters therefore shows that the total conditional expectation does not increase.

After all \(Q\) color stages every variable is fixed.
The objective is now a nonnegative integer still strictly below one, and hence equals zero.
There are \(Q\) stages of \(O(L(d+1))\) rounds plus the initial decomposition. \(\square\)

For our LCL, locate each red variable at its component leader, and \(E_K\) at the leader of \(K\).
A blue indicator may depend on several red variables, as different verifiers assigned to that blue component can touch different red components.
This causes no difficulty.

If a verifier assigned to blue \(K\) reads a red vertex \(u\) in \(U\), then some blue vertex \(v\in K\) lies in the same scope and \(d_G(u,v)\le2\rho\).
Each component leader is within \(D\) of each of its vertices.
The two leaders are consequently at distance at most \(2D+2\rho\).
Recovering the complete red components and their marked radius-\(T\) views requires a further \(O(D+s+T+1)\) radius, by Lemma 1 and the ambient-distance gathering argument.
Verifier descriptions, boundary inputs, and structural marks can be recovered within the same order of radius.
A sufficient common choice is

\[ L=20(D+s+T+\rho+1). \]

Thus the function describing \(E_K\), its variable locations, and rational tables from Lemma 6 are locally recoverable as required by Lemma 7.
There is no bound on the cardinality of a component, and none is needed with unbounded message size and local computation.

Use Lemma 7 to fix all red patterns with \(F=0\).
Distribute these patterns.
Each blue leader now executes its finite feasibility search, chooses the first passing pattern, and distributes it.
Every blue indicator vanishes, so every search succeeds.
No verifier reads blue coordinates from two components.
Every assigned verifier therefore accepts, and the already solved ordinary connected components remain valid.
The combined output belongs to \(\Pi(G)\).

All conditioning in this construction concerns the newly introduced product distribution.
We never condition \(\mu_G\) on global success or on the chosen red outputs.
In particular, we do not assume independence of blue outputs after revealing red outputs.

\section{7. Round bound and translation theorem}\label{round-bound-and-translation-theorem}

Combining the construction above proves the following precise statement.

\textbf{Theorem 8.} Under the model, effectiveness, metric-control, whole-set locality, bounded-dependence, and global-success hypotheses of Section 2, there is a uniform deterministic LOCAL algorithm solving \(\Pi\) in

\[ (T+s)(2+\log_2N)^{O(\log(C+2))}, \qquad s=1+\max\{k,2\rho,A,1\}, \]

rounds.
It requires neither a probability oracle nor a supplied metric decomposition.
The success assumption includes equality at \(1-1/n\).

\textbf{Proof of the round bound.} Lemma 4 costs \(O(r(P+d+1))\), with \(r=M^hs\).
Lemma 1 and input gathering recover the structural components and relevant marked views in \(O(D+s+T+\rho+1)\) rounds.
Table synthesis is a terminating local computation using no rounds.
Lemma 7 and final completion cost

\[ O\bigl((D+s+T+\rho+1)[P+Q(d+1)]\bigr). \]

Since \(D=Ms\le r\), \(s>\rho\), \(Q,d,P\) are polylogarithmic, and \(M=(4C+8)^4\), the total is the displayed bound.
In particular, factors polynomial in \(C\) can be absorbed into the exponent \(O(\log(C+2))\) because \(2+\log_2N\ge2\). \(\square\)

For a bounded-dependent outcome of locality \(T\), set \(k=2T\).
Then \(s=O(T+A+\rho+1)\) and

\[ \operatorname{DetLOCAL}(\Pi) \le (T+A+\rho+1)(2+\log_2N)^{O(\log(C+2))}. \]

Uniform Assouad--Nagata dimension at most one has fixed \(C\) and \(A=0\).
Fixed affine control also suffices.
Since \(n\le N\le n^c\), the factors \(2+\log_2N\) and \(2+\log_2n\) differ only by a constant factor.
This gives

\[ \operatorname{DetLOCAL}(\Pi) =O\bigl((T+\rho+1)\operatorname{polylog}n\bigr) \]

for fixed control constants.
The resulting algorithm finds a valid output; it need not preserve the original law or even belong to its support.

\section{8. Exact-support variant}\label{exact-support-variant}

There is also a probability-free version.
Suppose for each \(G\in\mathcal I_N\) there is a nonempty set \(S_G\subseteq\Pi(G)\) whose projection

\[ S_{G,U}=\{x|_U:x\in S_G\} \]

is determined by the marked radius-\(T\) input view, and which factorizes on arbitrary separated sets:

\[ d_G(U,V)>k\quad\Longrightarrow\quad S_{G,U\cup V}=S_{G,U}\times S_{G,V}. \]

Then a deterministic algorithm solves \(\Pi\) in

\[ O(T)+s(2+\log_2N)^{O(\log(C+2))} \]

rounds, with the same \(s\) as Theorem 8.

\textbf{Proof.} First synthesize a finite support table.
For each listed instance compute its finite set \(\Pi(G)\) by testing all output labelings.
Enumerate all choices of a nonempty subset of \(\Pi(G)\) for every instance.
Test each candidate for factorization on every pair of separated sets and for projection equality on every pair of isomorphic marked radius-\(T\) input views.
All tests are finite; the assumed support family is a passing candidate.
The first passing table is therefore computable from the common parameters and instance list without communication.
A component leader can recover its projection from its marked view, since tested locality makes all matching occurrences agree.

Construct the structural coloring at scale \(s\).
Each red leader gathers its marked radius-\(T\) view and chooses any pattern from its projected support.
Applying factorization to a red component and the union of all previous red components shows that all these choices jointly belong to \(S_{G,\text{red}}\).
Thus some \(x\in S_G\subseteq\Pi(G)\) agrees with every chosen red label.

Every red-only verifier already accepts.
Assign each verifier involving blue vertices to its unique blue component.
That component's leader gathers its assigned constraints and fixed red boundary and searches for a passing labeling.
The restriction of \(x\) witnesses that the search succeeds.
Independent blue completions satisfy every verifier, as in Section 6.
Gathering the red views and blue constraints costs \(O(D+s+T+\rho+1)\) beyond the structural construction; only the possible \(O(T)\) term is not covered by its scale. \(\square\)

A zero-error distribution satisfying the whole-set locality and dependence hypotheses provides such supports.
With positive error, taking the full support may include invalid outputs; the probability-free argument cannot simply be applied to it.
Conditioning the outcome on success also need not preserve independence or locality.
Sections 4--6 avoid both problems.
\clearpage
\part{Part II: The XOR-face separation on square grids without compass directions}
\section{1. The problem and the result}\label{grid-problem-and-result}

Let \(n\ge1\) be a perfect square and let
\[ G_n=P_{\sqrt n}\mathbin{\square}P_{\sqrt n}. \]
Here \(P_{\sqrt n}\) is the path on \(\sqrt n\) vertices, so \(G_n\) has \(n\) vertices and a boundary.
The input graph has no compass directions and no coordinate labels.
For the proof, we may draw its vertices at \((x,y)\in\{0,\ldots,\sqrt n-1\}^2\).
These coordinates are not given to the algorithm.

Each vertex knows its degree and has local port numbers \(1,\ldots,\deg(v)\) on its incident edges.
The port numbers at the two ends of an edge may be different.
Vertices have distinct identifiers in \([n^C]\), for a fixed constant \(C\ge1\), and may know the exact value of \(n\).
Correctness is required for every choice of identifiers and port numbers.
Messages and local computation are unrestricted, as in Part I.
We use input views that retain the original degrees and port numbers; another convention may change the round bounds by an additive constant.

Each vertex outputs four bits, one for each possible port.
A slot without an incident edge must contain zero.
For an edge \(e=uv\), write \(b_{u,e}\) and \(b_{v,e}\) for the bits output at its two ends.
A valid output must satisfy the following conditions:
\begin{enumerate}
\item The endpoints agree: \(b_{u,e}=b_{v,e}\) for every edge \(e=uv\).
Write \(X_e\) for their common value.
\item For every four-cycle \(Q\), the XOR of its four edge bits is one:
\begin{equation} \sum_{e\in E(Q)}X_e=1 \quad\text{in }\mathbb F_2. \label{eq:grid-face} \end{equation}
\end{enumerate}
In a square grid, every four-cycle is the boundary of a unit square.
We call this the \textbf{XOR-face problem}.
A vertex can check endpoint agreement on its incident edges within radius one, and it can check every four-cycle containing it within radius two.
Thus the problem is an LCL with verification radius two and output alphabet \(\{0,1\}^4\).

\begin{figure}[ht]
\centering
\setlength{\unitlength}{1pt}
\begin{picture}(245,130)
\put(55,25){\line(1,0){125}}
\put(55,25){\line(0,1){80}}
\put(55,105){\line(1,0){125}}
\put(180,25){\line(0,1){80}}
\put(55,25){\circle*{4}}
\put(180,25){\circle*{4}}
\put(180,105){\circle*{4}}
\put(55,105){\circle*{4}}
\put(45,13){\makebox(0,0){\(v_0\)}}
\put(190,13){\makebox(0,0){\(v_1\)}}
\put(190,117){\makebox(0,0){\(v_2\)}}
\put(45,117){\makebox(0,0){\(v_3\)}}
\put(118,13){\makebox(0,0){\(X_{e_1}\)}}
\put(202,65){\makebox(0,0){\(X_{e_2}\)}}
\put(118,117){\makebox(0,0){\(X_{e_3}\)}}
\put(33,65){\makebox(0,0){\(X_{e_4}\)}}
\end{picture}
\caption{The two endpoints store the same bit on each edge, and \(X_{e_1}+X_{e_2}+X_{e_3}+X_{e_4}=1\).
The drawing does not add directions or coordinates to the input.}
\label{fig:grid-face}
\end{figure}

An outcome assigns a joint distribution of outputs to each input graph.
It is \textbf{non-signaling with locality \(s\) on grids} if matching radius-\(s\) input views around two marked sets give the same joint distribution of all marked outputs.
The matching must preserve the marked vertices, graph edges, identifiers, port numbers, degrees, and supplied size information.
The marked sets may be large or disconnected.
Only comparisons between square-grid inputs are required here.

The outcome has \textbf{dependence range two} if output vectors on any finite collection of vertex sets at pairwise distance greater than two are mutually independent.
These are the locality and independence conditions of \href{https://arxiv.org/html/2504.05191v2\#S4}{Balliu et al., Definitions 4.2--4.4}, with the comparisons restricted to grids.
In the terminology of Part I, an outcome with these bounds is bounded-dependent with locality three, since independence beyond distance two also implies independence beyond distance six.

\textbf{Theorem (grid separation).}
On square grids without compass directions, the XOR-face problem has a zero-error outcome with non-signaling locality at most three and dependence range two.
Its deterministic LOCAL complexity is \(\Theta(\sqrt n)\), even with exact size information and unique identifiers.
The asymptotic statement is over perfect-square values of \(n\).

The rest of Part II proves all parts of this statement.
If \(n=1\), the only output is four zero bits, so the outcome claims are immediate.
The geometric arguments below concern grids with \(\sqrt n\ge2\).

\section{2. A distribution over valid outputs}\label{grid-distribution}

Let \(\Omega(G)\) be the set of edge-bit assignments satisfying \eqref{eq:grid-face} on every unit square of \(G\).
We define the outcome to be the uniform distribution on \(\Omega(G)\), with each edge bit copied to its two endpoint slots and all unused slots set to zero.
This definition uses only the graph and its four-cycles.
It does not choose a compass orientation.
We now prove that \(\Omega(G)\) is nonempty and describe this uniform distribution using independent vertex bits.

Choose coordinates only for this proof.
Define a fixed edge labeling \(\eta\) by
\begin{equation} \eta_{uv}=\begin{cases}0,&uv\text{ is horizontal},\\x\bmod2,&uv\text{ is vertical in column }x.\end{cases} \label{eq:grid-reference} \end{equation}
Around a unit square, the two horizontal bits are zero and the two vertical bits are the parities of consecutive columns.
Their XOR is one, so \(\eta\in\Omega(G)\).

Let \(Z_v\), one for each vertex, be independent uniform bits and set
\begin{equation} X_{uv}=\eta_{uv}+Z_u+Z_v. \label{eq:grid-random-bits} \end{equation}
All additions are modulo two.
Around a unit square, each vertex bit occurs on two incident edges and therefore cancels.
The remaining sum is the sum of the four \(\eta\)-bits, which is one.
Thus every sample is valid.

\textbf{Claim.}
Formula \eqref{eq:grid-random-bits} gives the uniform distribution on \(\Omega(G)\).

\textbf{Proof.}
Fix a valid edge labeling \(X\) and write \(a_e=X_e+\eta_e\).
The XOR of \(a_e\) around each unit square is zero.
Consider a simple cycle in the grid drawing and add these equations over all unit squares inside it.
Every interior edge appears twice and cancels, leaving XOR zero on the cycle boundary.
Every closed walk can be split into simple cycles and pairs of oppositely traversed edges, so it also has XOR zero.

Choose a root \(o\), and let \(z_v\) be the XOR of \(a_e\) along a path from \(o\) to \(v\).
Two such paths form a closed walk, so they give the same value.
For every edge \(uv\), this definition gives \(z_u+z_v=a_{uv}\).
Consequently, every valid \(X\) has the form \(\eta_{uv}+z_u+z_v\).

If two vertex-bit assignments give the same edge labeling, their difference has equal values at adjacent vertices.
The grid is connected, so their difference is constant on all vertices.
Hence each valid edge labeling has exactly two vertex-bit assignments producing it, and those two assignments are complements.
Independent uniform vertex bits give every one of the \(2^n\) assignments the same probability.
They therefore give every valid edge labeling the same probability, as claimed.
\(\square\)

Although \eqref{eq:grid-random-bits} uses coordinates in its proof, its output distribution is the coordinate-free uniform distribution just defined.
In particular, changing the coordinate origin, rotating the grid, or reflecting it leaves that distribution unchanged.
This does not give a fast sampling algorithm: the reference labeling \(\eta\) is not supplied to the vertices.

\section{3. Independence beyond distance two}\label{grid-independence}

Fix one input graph and one reference labeling \(\eta\).
The outputs at a set \(U\) depend only on the random bits at vertices in \(N_1(U)\), because every stored edge bit uses only its two endpoint bits.
If two marked sets are at distance greater than two, their radius-one neighborhoods are disjoint.
The same holds for any finite collection of sets at pairwise distance greater than two.
Their output vectors use disjoint collections of independent vertex bits and are therefore mutually independent.
This proves dependence range two.

Independence within one input does not by itself prove non-signaling between different inputs.
That comparison is proved next.

\section{4. Non-signaling for arbitrary marked sets}\label{grid-nonsignaling}

Let \(U\subseteq V(G)\) and \(U'\subseteq V(G')\) be marked sets in two square-grid inputs.
Suppose \(\varphi\) is an isomorphism between their radius-three input views, preserving all the input data listed in Section 1.
We prove that the entire output vector on \(U\) has the same distribution as the corresponding output vector on \(U'\).
This includes sets whose vertices are far apart.

Let \(F\) be the set of all edges incident to vertices of \(U\), and let \(W\) be their endpoints.
The marked outputs consist of the bits on these edges, with repetitions when both endpoints are marked, together with fixed zero padding.
The view isomorphism maps \(F\) to the corresponding edge set \(F'\) in \(G'\).
Different connected components of the graph \((W,F)\) have disjoint endpoint sets.
We will compare the edge labels on each such component jointly.

\subsection{4.1. What the graph tells us about directions}

For \(v\in W\), the radius-three view of \(U\) contains the full radius-two neighborhood of \(v\), since \(v\) is at distance at most one from \(U\).
This neighborhood determines whether two edges at \(v\) are perpendicular or opposite in the grid drawing.
Indeed, two distinct neighbors of \(v\) have a common neighbor other than \(v\) exactly when they lie in perpendicular directions.
That common neighbor completes a unit square.
Neighbors in opposite directions have no such second common neighbor.
These statements also hold at the boundary of a full square grid.

Use coordinates in both graph drawings for the proof.
At each \(v\in W\), the isomorphism therefore maps its incident coordinate directions by a rotation or reflection.
Write \(S_v\) for the corresponding linear map of the two coordinate axes.
It may exchange the axes or reverse their signs.
This map is unique: every vertex has two perpendicular incident edges, and the images of those two directions determine it.
The perpendicular and opposite relations then determine the images of all other incident directions.

We claim that \(S_v=S_w\) whenever \(vw\in F\).
The edge \(vw\) lies on at least one unit square, which we write in cyclic order as \(v,w,b,a\).
The whole square is visible, and its image is another unit square.
Let \(d=\chi(w)-\chi(v)\) be the coordinate direction from \(v\) to \(w\).
At \(v\), its image direction is \(S_vd\).
At \(w\), the reverse edge has direction \(-d\) and image direction \(S_w(-d)=-S_wd\), so the direction from \(\varphi(v)\) to \(\varphi(w)\) is also \(S_wd\).
Thus the two maps agree on \(d\).
They also agree on the perpendicular direction: the directions from \(v\) to \(a\) and from \(w\) to \(b\) are equal, and their image directions are equal.
They agree on two perpendicular directions and hence are the same linear map.

It follows that, on each connected component \(K\) of \((W,F)\), there is one rotation or reflection \(S_K\) and one integer translation \(c_K\) such that
\begin{equation} \chi'(\varphi(v))=S_K\chi(v)+c_K \qquad(v\in K), \label{eq:grid-coordinate-map} \end{equation}
where \(\chi\) and \(\chi'\) are the coordinate maps in the two drawings.
To see the translation, choose one vertex of \(K\) and then follow paths in \(K\).
On every edge the coordinate difference is transformed by the same \(S_K\), so \(\chi'(\varphi(v))-S_K\chi(v)\) is constant on the component.
Thus each observed-edge component is translated, rotated, or reflected as a whole, even though the input contains no compass directions.

\subsection{4.2. Matching all observed edge bits}

Choose reference labelings \(\eta\) and \(\eta'\) by \eqref{eq:grid-reference} in the two drawings.
Fix one component \(K\), write \(c_K=(a,b)\), and put \(\epsilon=a\bmod2\).
In the following bit calculations, coordinates are reduced modulo two.
We will choose fixed bits \(g_v\) so that
\begin{equation} \eta_{uv}+\eta'_{\varphi(u)\varphi(v)}=g_u+g_v \qquad(uv\in F\text{ in }K). \label{eq:grid-reference-difference} \end{equation}
Signs can be omitted in these bit calculations because \(-x=x\) modulo two.

If \(S_K\) does not exchange the axes, the two reference labelings differ by zero on horizontal edges and by \(\epsilon\) on vertical edges.
A vertical edge in column \(x\) becomes vertical in column \(\pm x+a\), so its two reference bits differ by \(x+(\pm x+a)=\epsilon\) modulo two.
Set \(g_v=(\epsilon y(v))\bmod2\).
Its endpoint sum is zero on a horizontal edge and \(\epsilon\) on a vertical edge, so \eqref{eq:grid-reference-difference} holds.

If \(S_K\) exchanges the axes, an original horizontal edge at row \(y\) becomes vertical in column \(\pm y+a\), while an original vertical edge becomes horizontal.
The reference-label difference is therefore \(y+\epsilon\) on a horizontal edge and \(x\) on a vertical edge.
Set \(g_v=(x(v)y(v)+\epsilon x(v))\bmod2\).
Across a horizontal edge its endpoint sum is \(y+\epsilon\), and across a vertical edge it is \(x\).
Again \eqref{eq:grid-reference-difference} holds.

Now generate the two output distributions using random bits related by
\begin{equation} Z'_{\varphi(v)}=Z_v+g_v \qquad(v\in W). \label{eq:grid-matched-random-bits} \end{equation}
Adding fixed bits to independent uniform bits leaves them independent and uniform.
The components have disjoint endpoint sets, so we can make this choice on all components at once and choose the remaining primed vertex bits independently.
For every observed edge \(uv\), equations \eqref{eq:grid-reference-difference} and \eqref{eq:grid-matched-random-bits} give
\[ \eta'_{\varphi(u)\varphi(v)}+Z'_{\varphi(u)}+Z'_{\varphi(v)}=\eta_{uv}+Z_u+Z_v. \]
Thus corresponding observed edge labels are equal under this joint construction.
The view isomorphism preserves the port slots and the unused slots, so the full marked output vectors are equal as well.
Each side has the required uniform outcome distribution, by Section 2.
Their entire joint marginals therefore agree, proving non-signaling locality at most three.

\section{\texorpdfstring{5. Why deterministic algorithms need \(\Omega(\sqrt n)\) rounds}{5. Why deterministic algorithms need Omega(sqrt(n)) rounds}}\label{grid-lower-bound}

\subsection{5.1. The parity contradiction on an odd torus}

Let \(m\ge5\) be odd and let \(H=C_m\mathbin{\square}C_m\).
Think of \(H\) as an \(m\)-by-\(m\) grid with opposite sides joined.
It has \(m^2\) unit faces, and each edge belongs to exactly two of them.
Suppose an edge labeling satisfied XOR one around every face.
Adding the \(m^2\) equations modulo two would give zero on the left, since every edge bit appears twice, and one on the right, since \(m^2\) is odd.
This is impossible.
Thus no consistent assignment of edge bits satisfies all face equations on \(H\).

We will use this contradiction to prove the lower bound on square grids with boundary.
The random outcome is not evaluated on the torus.

\subsection{5.2. Comparing torus neighborhoods with square-grid inputs}

Suppose a deterministic algorithm \(A\) solves the LCL on every \(n\)-vertex square grid and every identifier and port assignment in \(t=t(n)\) rounds.
Let \(n\ge144\) be a perfect square and suppose, for a contradiction, that
\begin{equation} t+3\le\sqrt n/4. \label{eq:grid-fast-algorithm} \end{equation}
Let \(m\) be the largest odd integer at most \(\sqrt n\), so \(m\ge\sqrt n-1\) and \(m^2\le n\).
Use the torus \(H=C_m\mathbin{\square}C_m\), give it distinct identifiers from \([n]\), and choose any port numbering.
In all its local views, use the same public size parameter \(n\) as on the square-grid inputs.
The torus may have fewer than \(n\) vertices.
We do not assume that \(A\) is correct on it; we use only local views that occur in actual \(n\)-vertex square grids.

Put \(r=t+3\).
The assumed bound gives \(2r+1<m\).
A radius-\(r\) neighborhood in the torus therefore has exactly the same vertices and edges as a radius-\(r\) neighborhood in the square lattice.
In more detail, two vertices of such a lattice ball cannot be identified modulo \(m\), and an extra torus edge would require \(m\le2r+1\).

This neighborhood fits in the interior of \(G_n\).
Place its center at
\[ \left(\left\lfloor\frac{\sqrt n-1}{2}\right\rfloor,\left\lfloor\frac{\sqrt n-1}{2}\right\rfloor\right). \]
Every coordinate changes by at most \(r\le\sqrt n/4\) inside the ball.
For \(n\ge144\), this leaves room before the physical boundary of the square grid.
Copy all identifiers and port information needed by the local views, and give the remaining square-grid vertices distinct unused identifiers from \([n]\).
Choose the remaining port numbers arbitrarily.
This gives a genuine input with exactly \(n\) vertices.
The extra layers in \(r=t+3\) also preserve the degree and port information needed at the four vertices of a face.

Every radius-\(t\) view at a torus vertex is therefore a view that can occur at a vertex of a valid square-grid input with public size \(n\).
Apply the deterministic output rule of \(A\) to that view and call the resulting port tuple \(Y_v\).
Doing this for every torus vertex fixes one output assignment \(Y\) on \(H\).
This definition does not require \(A\) to be correct on a torus.

\subsection{5.3. Every face forces a constraint on that same output}

Fix a torus face \(Q\) and choose one of its vertices \(v\).
All four vertices of \(Q\) are within distance two of \(v\).
The radius-\(t\) views determining their outputs, including the needed boundary information, are contained in the radius-\(t+3\) view of \(v\).
Embed that view in an \(n\)-vertex square grid as above, preserving identifiers, ports, degrees, and the public size.
The four outputs produced there by \(A\) are exactly the four tuples \(Y_w\) already fixed on the torus face.

Correctness on this genuine input implies agreement of the endpoint bits on all four edges of \(Q\), and XOR one around \(Q\).
The argument applies to every torus face.
A different square-grid input may be used for each face, but the assignment \(Y\) was fixed before these comparisons.
Hence all the equations concern the same torus output assignment.
Every torus edge belongs to a face, so endpoint agreement gives a single bit \(X_e\) on every edge, and every face has XOR one.
This contradicts Section 5.1.

We conclude that \(t(n)>\sqrt n/4-3\) for every perfect-square \(n\ge144\), under the stated input-view convention.
In particular, the deterministic round complexity is \(\Omega(\sqrt n)\).
This argument preserves exact size information and numerical identifiers and uses no compass input.

\section{6. A matching deterministic upper bound}\label{grid-upper-bound}

The grid has diameter \(2(\sqrt n-1)\).
In \(O(\sqrt n)\) rounds, every vertex can gather the whole graph, including identifiers and port numbers.
Order the edges lexicographically by the pairs of their endpoint identifiers, putting the smaller identifier first, and enumerate all bit assignments in that order.
Choose the first assignment satisfying every face equation.
Such an assignment exists by Section 2.
Every vertex sees the same complete input and uses the same ordering, so they choose the same assignment and copy its edge bits to their incident ports.
The search is finite, and the LOCAL model does not charge for its computation time.
This gives a deterministic \(O(\sqrt n)\)-round algorithm and proves the claimed \(\Theta(\sqrt n)\) complexity.

\section{7. Graph families and the translation theorem}\label{grid-relation-to-part-one}

The example shows that a distribution with constant non-signaling locality, constant dependence range, and zero error can still require \(\Theta(\sqrt n)\) rounds for deterministic solutions when the non-signaling comparisons are restricted to grids.
Square grids have Assouad--Nagata dimension two, so this rules out a general dimension-two extension based only on those grid comparisons.
Part I requires dimension-one metric control and makes no such dimension-two claim.

There is also a reason why the grid outcome cannot simply be extended to all bounded-degree graphs while retaining a fixed non-signaling radius.
Take a sufficiently large odd integer \(\sqrt n\) and the torus \(C_{\sqrt n}\mathbin{\square}C_{\sqrt n}\), which has exactly \(n\) vertices.
For any fixed radius, the neighborhoods needed by each face and its edge-consistency tests fit inside a genuine \(n\)-vertex square grid, with the same identifiers, ports, and public size.
An all-graphs non-signaling outcome that is always correct on square grids would make each of these torus tests pass with probability one.
There are finitely many tests, so they would all pass simultaneously with probability one.
Section 5.1 shows that this is impossible.
This separate argument uses only non-signaling; finite dependence is not needed for it.

The result proved here concerns full square grids with boundary and all assignments of identifiers and ports.
We do not extend the no-compass statement to arbitrary lattice subgraphs with holes or deleted edges.
The LCL framework is due to Naor and Stockmeyer, and the grid setting is studied by Brandt et al.; the references below give these sources.
The proofs in Part II are complete and do not use a grid classification theorem.

\clearpage
\section{References}
\begin{enumerate}
\item Alkida Balliu, Filippo Casagrande, Francesco d'Amore, Massimo Equi, Barbara Keller, Henrik Lievonen, Dennis Olivetti, Gustav Schmid, and Jukka Suomela.
\emph{Distributed Quantum Advantage in Locally Checkable Labeling Problems}.
SODA 2026, pp. 1268--1308.
Expanded version \href{https://arxiv.org/abs/2504.05191v2}{arXiv:2504.05191v2}, 2026.
Definition and theorem numbers refer to this version.
\item Francesco d'Amore and Henrik Lievonen.
\emph{Superlogarithmic Gap Result for LCLs on Trees in Quantum-LOCAL}.
\href{https://arxiv.org/abs/2608.16854v1}{arXiv:2608.16854v1}, 2026.
\item Marc Distel.
\emph{Proper Minor-Closed Classes of Graphs Have Assouad--Nagata Dimension 2}.
\href{https://arxiv.org/abs/2308.10377v3}{arXiv:2308.10377v3}, 2025.
Cited numbering refers to this version.
\item Chun-Hung Liu.
\emph{Assouad--Nagata dimension of minor-closed metrics}.
Proceedings of the London Mathematical Society 130 (2025), e70032.
\href{https://arxiv.org/abs/2308.12273v5}{arXiv:2308.12273v5}; \href{https://doi.org/10.1112/plms.70032}{doi:10.1112/plms.70032}.
\item Mohsen Ghaffari and Christoph Grunau.
\emph{Near-Optimal Deterministic Network Decomposition and Ruling Set, and Improved MIS}.
FOCS 2024.
\href{https://arxiv.org/abs/2410.19516v1}{arXiv:2410.19516v1}.
\item V\'aclav Rozho\v{n} and Mohsen Ghaffari.
\emph{Polylogarithmic-Time Deterministic Network Decomposition and Distributed Derandomization}.
STOC 2020.
\href{https://arxiv.org/abs/1907.10937}{arXiv:1907.10937}.
\item Moni Naor and Larry Stockmeyer.
\emph{What Can Be Computed Locally?} SIAM Journal on Computing 24(6):1259--1277, 1995.
\href{https://doi.org/10.1137/S0097539793254571}{doi:10.1137/S0097539793254571}.
\item Sebastian Brandt, Juho Hirvonen, Janne H. Korhonen, Tuomo Lempi\"ainen, Patric R. J. \"Osterg\aa rd, Christopher Purcell, Joel Rybicki, Jukka Suomela, and Przemys\l aw Uzna\'nski.
\emph{LCL problems on grids}.
PODC 2017.
\href{https://arxiv.org/abs/1702.05456v2}{arXiv:1702.05456v2}.
\end{enumerate}

\end{document}
